\section{Đối chiếu với implementation trong repo}
\begin{frame}{Paper \texorpdfstring{$\leftrightarrow$}{<->} code repo}
  \begin{table}\centering\tiny
    \begin{tabular}{p{2.9cm}p{6.9cm}}\toprule
      \textbf{Khối trong paper} & \textbf{File trong repo}\\\midrule
      Item image summary & \texttt{Inference/microlens/}\texttt{image\_summary.py} — tạo mô tả ảnh bằng LLaVA; đã verify file tồn tại.\\
      Recurrent preference & \texttt{Inference/microlens/}\texttt{preferece\_inference\_recurrent.py} — chunk 3 item, cập nhật summary bằng Llama 3; đã verify file tồn tại.\\
      Dataset construction & \texttt{train/microlens/}\texttt{dataset\_create.py}, \texttt{test/microlens/}\texttt{multi\_col\_dataset.py}; đã verify file tồn tại.\\
      MLLM SFT & \texttt{train/microlens/}\texttt{train\_llava\_sft.py} — LLaVA + LoRA + PyTorch Lightning; đã verify file tồn tại.\\
      Evaluation & \texttt{test/microlens/}\texttt{test\_with\_llava\_sft.py}; đã verify file tồn tại.\\\bottomrule
    \end{tabular}
  \end{table}
  \vspace{0.15cm}\scriptsize\textcolor{Logo1}{Lưu ý:} đây là đối chiếu implementation trong repo, không thay thế thông số được báo cáo trong paper.
\end{frame}

\begin{frame}{Một số sai khác cần ghi nhớ khi tái hiện}
  \begin{itemize}
    \item \textbf{Đã verify trực tiếp tại `train/microlens/train\_llava\_sft.py`:} paper báo cáo LoRA rank 8, 10 epochs và max length 512; checkout hiện tại đặt `LORA\_R=16`, `EPOCH=4`, `MAX\_LENGTH=1024` (dòng 1--4).
    \item \textbf{Đã verify trực tiếp trong checkout hiện tại:} paper dùng Amazon-Baby/Amazon-Game; repo có các file xử lý `data/amazon/processed/Baby\_Products\_*` và `Video\_Games\_*`, không dùng đúng tên thư mục dataset của paper.
    \item `accumulate\_grad\_batches=4` trong repo hiện tại (dòng 242--248), trong khi paper báo cáo gradient accumulation 8; đây là khác biệt cấu hình cần giữ riêng.
    \item README có một số đường dẫn script cũ; khi chạy cần dùng path thực tế dưới \texttt{MLLM-MSR/train}, \texttt{MLLM-MSR/test} và \texttt{Inference/microlens}.
  \end{itemize}
  \begin{block}{Bài học tái lập}
    Phải phân biệt “con số trong paper” với “cấu hình của code checkout hiện tại”; nếu không, kết quả chạy lại không nhất thiết khớp bảng báo cáo.
  \end{block}
\end{frame}
